% Options for packages loaded elsewhere
\PassOptionsToPackage{unicode}{hyperref}
\PassOptionsToPackage{hyphens}{url}
\PassOptionsToPackage{dvipsnames,svgnames,x11names}{xcolor}
%
\documentclass[
  11pt,
  a4paper,
]{ltjsarticle}
\usepackage{amsmath,amssymb}
\usepackage{iftex}
\ifPDFTeX
  \usepackage[T1]{fontenc}
  \usepackage[utf8]{inputenc}
  \usepackage{textcomp} % provide euro and other symbols
\else % if luatex or xetex
  \usepackage{unicode-math} % this also loads fontspec
  \defaultfontfeatures{Scale=MatchLowercase}
  \defaultfontfeatures[\rmfamily]{Ligatures=TeX,Scale=1}
\fi
\usepackage{lmodern}
\ifPDFTeX\else
  % xetex/luatex font selection
\fi
% Use upquote if available, for straight quotes in verbatim environments
\IfFileExists{upquote.sty}{\usepackage{upquote}}{}
\IfFileExists{microtype.sty}{% use microtype if available
  \usepackage[]{microtype}
  \UseMicrotypeSet[protrusion]{basicmath} % disable protrusion for tt fonts
}{}
\makeatletter
\@ifundefined{KOMAClassName}{% if non-KOMA class
  \IfFileExists{parskip.sty}{%
    \usepackage{parskip}
  }{% else
    \setlength{\parindent}{0pt}
    \setlength{\parskip}{6pt plus 2pt minus 1pt}}
}{% if KOMA class
  \KOMAoptions{parskip=half}}
\makeatother
\usepackage{xcolor}
\usepackage[margin=24mm]{geometry}
\setlength{\emergencystretch}{3em} % prevent overfull lines
\providecommand{\tightlist}{%
  \setlength{\itemsep}{0pt}\setlength{\parskip}{0pt}}
\setcounter{secnumdepth}{5}
\ifLuaTeX
  \usepackage{selnolig}  % disable illegal ligatures
\fi
\IfFileExists{bookmark.sty}{\usepackage{bookmark}}{\usepackage{hyperref}}
\IfFileExists{xurl.sty}{\usepackage{xurl}}{} % add URL line breaks if available
\urlstyle{same}
\hypersetup{
  colorlinks=true,
  linkcolor={blue},
  filecolor={Maroon},
  citecolor={Blue},
  urlcolor={blue},
  pdfcreator={LaTeX via pandoc}}

\author{}
\date{}

\begin{document}

\hypertarget{ux89e3ux7b54-11-ux6570ux5024ux7ddaux5f62ux4ee3ux6570}{%
\section{解答 11 ---
数値線形代数}\label{ux89e3ux7b54-11-ux6570ux5024ux7ddaux5f62ux4ee3ux6570}}

\hypertarget{section}{%
\subsection{1}\label{section}}

特異値は \texttt{1,10\^{}\{-8\}} なので条件数は \texttt{10\^{}8}
です。単位円が一方向に \texttt{10\^{}\{-8\}}
まで潰れた極端に細い楕円へ写ります。

\hypertarget{section-1}{%
\subsection{2}\label{section-1}}

\texttt{A\^{}TA=diag(1,10\^{}\{-16\})} なので条件数は
\texttt{10\^{}16=κ(A)\^{}2}。有限精度では小さい方向を扱う有効桁がさらに失われやすくなります。

\hypertarget{section-2}{%
\subsection{3}\label{section-2}}

悪条件性は入力の微小変化に真の解が敏感な性質です。不安定性は、well-conditioned
な問題でも計算法自体が不要な誤差を増幅する性質です。良いアルゴリズムでも悪条件問題の感度そのものは消せません。

\hypertarget{section-3}{%
\subsection{4}\label{section-3}}

疎性を保てばメモリと計算量を大幅に削減できます。逆行列は一般に dense
になり得ますが、\texttt{Ax} は疎構造をそのまま利用できます。

\hypertarget{section-4}{%
\subsection{5}\label{section-4}}

完全 SVD は不要な小特異成分まで計算するため高コストです。randomized SVD
はランダムベクトルを \texttt{A}
に作用させ、像の主要部分空間を低次元で捕捉してから、その小さな空間内で
SVD を行います。

\end{document}
